% Nouveau EGA V (cote 25): archive cote 25, site page 118 (PDF page 119).
% typescript. Draft compared: none.
% AI-transcribed by Claude Opus 5.5 (claude-opus-5-5, adaptive thinking, effort high), Message Batches API, batch ega25b (msgbatch_01A7n5YuXubUPNK1zsepi8RG), 2026-11-13; not yet checked by eye.
\scanpage{25/118}
\mspage{\hand{21}}

ce qui entrainera la conclusion voulue. Soit $\underline{J}$ l'Idéal qui définit
D , il suffira évidemment de prouver que $\underline{J}\to i_{\hand{*}}(\underline{J}|X-Z")$ est un isomor\xout{p}
phisme ($i:X-Z"\to X$ désignant l'immersion canonique), alors l'homomorphis
me $\underline{J}\to\underline{O}_X$ se reconstruit en effet en appliquant le foncteur $i_{\hand{*}}$ à
$\underline{J}\,|\,X-Z" \longrightarrow \underline{O}_X\,|\,X-Z"$ . Or comme $\underline{J}$ est inversible, il est plat sur S \unsure{et}
\xout{fibresxde} $x\in Z" \Rightarrow \mathrm{prof}_x\,\underline{J}_{s\,\hand{\unsure{*}}} \geqslant 2$. Il suffit donc de prouver le
\uline{Lemme} \hand{\unsure{3.2|3.1}}. Soit $X\to S$ loc de présentation finie, F un Module de présentation
finie sur X plat rel à S, T une partie fermée de X, supposons que pour
\margin{\unsure{$\downarrow$ \S 20 !!}}
tout $x\in X$ au dessus de $s\in S$, on ait \hand{\unsure{/}} $\mathrm{prof}_x F_{s\,\hand{\unsure{*}}}\geqslant 1$ (resp. $\mathrm{prof}_x F_s \geqslant 2$).
Alors l\xout{e}'homomorphisme canonique $F\to i_{\hand{*}}(F|X-T)$ est injectif (resp.
bijectif), où $i:X-T\to T$ est l'immersion canonique.

Démonstration du lemme: on peut supposer S,X affines, et on se
ramène par le procédé bréveté au cas S noethérien. Alors l'hypothèse
implique, en vertu du par.6 , que l'on a $\mathrm{prof}_x F \geqslant \mathrm{prof}_x F_s$ pour tout .
$x\in X$ sur $s\in S$, donc $\mathrm{prof}_x F_s \geqslant 1$ (resp.$\geqslant 2$) si $x\in T$. On conclut alors
grâce au par.5 . (NB Pour bien faire, ce lemme devrait figurer au par.
11 sous la rubrique; élimination des hypothèses noethériennes ...).

Il reste enfin à vérifier la dernière assertion\xout{s} de prop.2, savoir
que \xout{\unsure{Z}}\inserted{X-Z} est exactement le domain\xout{e} de définition de l"application rationel-%
le rel à S définie par f. Soit $U\supset X-Z$ ce domaine, il résulte alors de
la proposition 1 que $U\longrightarrow P$ est associé à un\xout{xfamillex}\inserted[typed]{système lin.} de diviseurs
D' sur U\unsure{/}S paramétré par $\check{P}$ , et on a $D'|(X-Z)\times_S\check{P} = D|(X-Z)\times_S\check{P}$ . Appli
quant le résultat d'unicité \hand{(}déjà démontré\hand{)}à D' et $D|U\times_S\check{P}$ , on voit donc
que ces deux derniers diviseurs sont égaux, donc $D|U\times_S\check{P}$ n'a pas de
point fixes i.e. $Z\cap U=\emptyset$ donc U=X-Z , cqfd.

\mspage{\hand{118}}
